% The ECG viewer in full: the formats it reads (WFDB single- and multi-segment, GE MUSE XML,
% DICOM ECGs, EDF and BDF, its JSON form, PNG, JPEG, PDF), opening one ECG or a folder, the
% paper, the calipers, the report, the themes, the addresses, muse-py (serve, json, wfdb, index),
% the examples.
% Sources: viewers/ecg/README.md, viewers/ecg/python/muse/README.md, viewers/ecg/examples/README.md,
% docs/edf.md, checked against viewers/ecg/index.html, app.js, readers.js (the registry and the
% record), wfdb.js, muse.js, dicom-ecg.js, edf.js, ecg.js, paper.js, style.css and
% viewers/ecg/python/muse/muse_reader_py/cli.py and page_server.py.
\chapter{The ECG viewer}
\label{chap:ecg-viewer}

This chapter describes the ECG viewer: the files it reads, how to open them, the paper and its
tools, and how to run it alone or with \cmd{muse-py}. It is written for clinicians and
researchers who read ECGs, and for administrators who deploy the viewer.

The ECG viewer shows 12-lead ECGs as cardiologists read them: on standard ECG paper, with their
report. It is one of Medorama's viewers, and also a page of its own that needs no installation.
Open it on the website (\url{https://medorama.org/viewers/ecg/}), or open
\file{viewers/ecg/index.html} from a clone, from disk or from any web server; then choose
\ui{Open record\dots} or \ui{Open multiple records\dots}, or drop the files on the page. The
files are read and drawn in the browser, and sent nowhere. The page has a light and a dark theme,
and follows the system's setting until one is chosen (\cref{sec:ecg-viewer-themes});
\cref{fig:ecg-viewer-main,fig:ecg-viewer-dark} show both.

\screenshot{ecg.png}{The ECG viewer in the light theme, as Medorama shows it (the bar above the paper is Medorama's, with a synthetic patient): a synthetic 12-lead ECG in the layout 3~$\times$~4 with a rhythm strip, and its report above the paper.}{fig:ecg-viewer-main}

\screenshot{ecg-dark.png}{The same ECG in the dark theme.}{fig:ecg-viewer-dark}

% ==========================================================================================
\section{What it reads}
\label{sec:ecg-viewer-formats}

\begin{xltabular}{\linewidth}{@{}P{4.2cm}L@{}}
\caption{The files the ECG viewer reads.}
\label{tab:ecg-viewer-formats}\\
\toprule
{\tabhead File} & {\tabhead Shown} \\
\midrule
\endfirsthead
\toprule
{\tabhead File} & {\tabhead Shown} \\
\midrule
\endhead
\bottomrule
\endfoot
WFDB record: the \file{.hea} header with its signal file(s), chosen together & Drawn on paper,
  in the signal formats 8, 16, 24, 32, 61, 80, 160 and 212. The first 10\,s are read, and the
  header's comments (\code{\# QT: 410 ms}, \code{\# Dx: \dots}) are the report. A skewed signal
  (\code{16:3}) is read that many frames later, as the WFDB library reads it. Where the header
  leaves a value out, WFDB's own defaults apply (\cref{sec:ecg-viewer-wfdb}). \\
WFDB multi-segment record: the master header with its segments' headers and signal files & The
  same. The segments follow each other from the record's start; a null segment, and a lead that a
  segment does not hold, are gaps. \\
GE MUSE XML (resting ECG) & Drawn on paper (the rhythm strip, with III, aVR, aVL and aVF
  derived), with its report: details, measurements, the Marquette 12SL statements and the
  confirmed or edited ones. \\
DICOM ECG (12-lead, Holter or other waveform object) & Drawn on paper, with its annotations as
  the report. Any other DICOM file is refused, with a pointer to the PACS viewer. \\
EDF, EDF+, BDF, BDF+ (\file{.edf}, \file{.bdf}) & Drawn on paper: the first 10\,s of the ECG
  signals, with the annotations of that span as notes (\cref{sec:ecg-viewer-edf}). \\
The JSON form (\cref{sec:ecg-viewer-json}) & Drawn on paper. \\
PNG, JPEG (and GIF, WebP) & Shown as it is (a scanned or exported ECG). The layout, the speed,
  the gain and the calipers are grayed out. \\
PDF & Shown as it is, in the browser's PDF viewer, with the same controls grayed out. \\
\end{xltabular}

A file is recognized by its extension, or else by its first bytes. A file of another kind is
refused with \code{This is not an ECG the viewer reads} and the list of what the viewer reads.
Each format has its own reader, and every reader produces the same \emph{record}
(\cref{sec:ecg-viewer-files}; adding a format is described in \cref{sec:readers-ecg}).

\subsection{WFDB records}
\label{sec:ecg-viewer-wfdb}

A WFDB record is its header (\file{.hea}) and the signal files the header names, which must be
chosen with it (\code{Choose the header together with its signal file FILE.} otherwise). The
viewer reads the first 10 seconds of every signal, in the formats 8, 16, 24, 32, 61, 80, 160 and
212, converts them to millivolts with each signal's gain, baseline and units (\code{mV},
\code{uV} or \code{V}), and keeps a signal with several samples per frame at its own rate. The
header's start time is shown as UTC.

\begin{note}
Where the header leaves a value out, the ECG viewer applies WFDB's own defaults: a sampling
frequency of 250\,Hz; a gain of 200\,adu/mV, also for a gain of 0; the unit mV; the baseline at
the ADC zero. A header that states them all is read as stated. VitalStream, unlike this viewer,
reads only what a header states (\cref{sec:viewers-alone-opening}).
\end{note}

\paragraph{The report in the header's comments.} The comment lines (\code{\# Key: value}, or
\code{Key = value}) are the report above the paper:
\begin{itemize}
\item a measurement under one of its usual names, with a number, is a measurement, with its
  default unit when none is given: \code{Vent. rate} (or \code{HR}, \code{Heart rate}),
  \code{Atrial rate}, \code{PR}, \code{QRS}, \code{QT}, \code{QTc}, \code{QTcB}, \code{QTcF},
  \code{RR}, \code{P axis}, \code{R axis} (or \code{QRS axis}), \code{T axis};
\item a key that ends with a lead's name (\code{ST V5: -0.12 mV}) is a measurement of that lead;
\item \code{Dx}, \code{Diagnosis}, \code{Interpretation}, \code{Statement}, \code{Comment},
  \code{12SL} and \code{Reason} are the computer's statements; \code{Program} (or
  \code{12SL version}) names its program;
\item \code{Dx confirmed}, with \code{Confirmed by} and \code{Confirmed at}, are the physician's
  confirmed statements; \code{Dx edited} the edited ones, not confirmed. Statements in one line
  are separated by \code{;} or \code{|};
\item any other \code{key: value} is a detail of the ECG; a line without a key is a note;
\item a key that names the patient anywhere in it (\code{Name}, \code{Full name},
  \code{Patient ID}, \code{Subject ID}, \code{MRN}, \code{Medical record number},
  \code{Accession number}, \code{NHS number}, \code{Date of birth}, \code{Born}, \code{Address},
  \code{Zip}, \code{Phone}, \code{Email} and the like) is never shown, even when the file holds
  it. Free text, a comment line without a key and the statements, is shown as written.
\end{itemize}

\begin{lstlisting}[style=text]
rest_0001 12 500 5000 09:30:00 02/01/2024
rest_0001.dat 16 1000/mV 16 0 0 0 0 I
...
# Vent. rate: 72 bpm
# QTc: 412 ms
# ST V5: -0.12 mV
# Dx: Sinus rhythm; ST depression, inferolateral leads
\end{lstlisting}

\subsection{Multi-segment records}
\label{sec:ecg-viewer-multisegment}

A multi-segment record (a master header whose record line is \code{name/n}) is chosen with its
segments' headers and signal files. Its segments follow each other from the record's start; a
null segment (\code{\textasciitilde}) is a gap of its length in every lead, a lead that a segment
does not hold is a gap in that lead, and the layout segment of a variable layout names the leads.
The first 10 seconds are read, as for one record.

\subsection{GE MUSE XML}
\label{sec:ecg-viewer-muse}

The page reads GE MUSE resting ECG XML files (\code{RestingECG}) with its own small XML parser,
\file{muse.js}; another XML file is refused (\code{This XML file is not a MUSE resting ECG
(RestingECG).}). Its Python twin, \cmd{muse-py} (\cref{sec:ecg-viewer-muse-py}), gives the same
result; both are tested on the same synthetic file.

\begin{table}[htbp]
\centering
\caption{What is read from a MUSE XML file.}
\label{tab:ecg-viewer-muse}
\begin{tabularx}{\linewidth}{@{}P{2.6cm}L@{}}
\toprule
\tabhead Part & \tabhead From \\
\midrule
Leads & The rhythm waveform (\code{Waveform} of type \code{Rhythm}): base64 little-endian 16-bit
  integers times \code{LeadAmplitudeUnitsPerBit}. III, aVR, aVL and aVF are derived from I and
  II, as the carts derive them \\
Sampling rate & \code{SampleBase} $\times$ 10 to the power \code{SampleExponent}. A base that
  is not a positive number is taken as 500\,Hz, and an exponent beyond $\pm$6 is refused \\
Details & \code{TestDemographics}: the acquisition time, the device, the site, the location, the
  reason, the status, and the filters; the patient's age and sex \\
Measurements & \code{RestingECGMeasurements}: the ventricular and atrial rates, PR, QRS, QT, QTc,
  QTcF, the P, R and T axes, the QRS count \\
Statements & \code{OriginalDiagnosis} gives the computer's (Marquette 12SL, with its version).
  \code{Diagnosis} gives the physician's: \emph{confirmed} (with the overreader and the time of
  the edit) when the status is \code{CONFIRMED}; \emph{edited} when it is not confirmed but
  differs from the computer's \\
\bottomrule
\end{tabularx}
\end{table}

The patient's name, ID and date of birth, the order (accounts, admission, physicians' IDs) and
other identifiers are never read into what the page shows.

\subsection{DICOM ECGs}
\label{sec:ecg-viewer-dicom}

A DICOM 12-lead ECG, a Holter ECG, or another DICOM waveform object, is read by the PACS viewer's
DICOM reader, of which the ECG viewer keeps a copy (\file{dicom.js}). The viewer draws the longest
group of channels (the rhythm, not the median beats), in microvolts unless the file says
otherwise, and shows the file's annotations as the report: its measurements, under the names the
viewer uses (\code{QT}, \code{QRS}, \code{Vent. rate}, \dots), its text as the computer's
statements, and the acquisition time and the group's label as details. Any other DICOM file is
refused: \code{This DICOM file is not an ECG: open it in the PACS viewer.}

\subsection{EDF, EDF+, BDF and BDF+}
\label{sec:ecg-viewer-edf}

The viewer reads files of the European Data Format family with its own reader, \file{edf.js},
which follows the same reading rules as the engines of VitalStream and Medorama
(\cref{sec:formats-reference-edf}). Since the viewer draws a 12-lead-style ECG on paper, it reads
a short span of the file:
\begin{itemize}
\item \textbf{The leads} are the ordinary signals at the sampling rate of the first ECG signal:
  a signal labelled \code{I} to \code{V6}, or whose label starts with \code{ECG} or \code{EKG}.
  A file without such a signal uses the signals at the most common rate. The other signals are
  left out, and a note lists them. Samples are converted to millivolts.
\item \textbf{The span} is the first 10 seconds of the recording's first uninterrupted part. An
  interrupted (EDF+D) file is read up to its first break, with a note.
\item \textbf{Annotations} within that span are listed as notes, with their times from the start
  of the span: at most 100, followed by a count of the others.
\item \textbf{The header's identification fields} are never read, with one exception: the four
  digits of the year of an EDF+ recording field's \code{Startdate dd-MMM-yyyy}, which date the
  recording.
\item The page names the format (\code{EDF}, \code{EDF+}, \code{BDF} or \code{BDF+}), and the
  signals' filters when they share one.
\end{itemize}
Long recordings, such as sleep studies and EEGs, are best opened in VitalStream
(\cref{sec:vitalstream-edf}). In Medorama, an EDF file of the category \code{ecg} is an ECG
attachment that opens in this viewer; in Privacy Mode it is sent de-identified
(\cref{sec:protecting-edf}).

\subsection{The JSON form}
\label{sec:ecg-viewer-json}

The JSON form is a file that any program can write (a MUSE XML reader, for instance). It holds
the samples of each lead, in the lead's units, with \code{null} where a sample is missing. It is
the ECG record of \cref{sec:readers-ecg-record} written as JSON.

\begin{lstlisting}[language=JSON,style=example]
{
  "name": "ECG 2024-01-02 09:30",
  "start": "2024-01-02T09:30:00Z",
  "fs": 500,
  "leads": [
    { "name": "I", "units": "uV", "samples": [12, 15, 19] },
    { "name": "II", "units": "uV", "samples": [20, 24, 30] }
  ],
  "meta": {
    "fields": [["Device", "MAC 5500"], ["Filters", "0.05-150 Hz"]],
    "measurements": [{ "name": "QTc", "value": 412, "unit": "ms" },
                     { "name": "ST", "value": -0.07, "unit": "mV", "lead": "V5" }],
    "statements": [{ "source": "computer", "program": "12SL", "lines": ["Sinus rhythm"] }],
    "notes": ["synthetic"]
  }
}
\end{lstlisting}

\begin{xltabular}{\linewidth}{@{}P{3.4cm}L@{}}
\caption{The fields of the ECG viewer's JSON form.}
\label{tab:ecg-viewer-json}\\
\toprule
{\tabhead Field} & {\tabhead Meaning} \\
\midrule
\endfirsthead
\toprule
{\tabhead Field} & {\tabhead Meaning} \\
\midrule
\endhead
\bottomrule
\endfoot
\code{name} & The ECG's name; default \code{ECG} \\
\code{start} & The time of the first sample: an ISO 8601 text, or Unix seconds; optional \\
\code{fs} & The sampling rate, in Hz; it may instead be given per lead \\
\code{leads} & Required: one object per lead \\
\quad\code{name} & The lead's name (\code{I}, \code{aVR}, \code{V1}, \code{ECG II}, \dots) \\
\quad\code{units} & \code{mV}, \code{uV} (or \code{$\mu$V}) or \code{V}; other units are taken as
  millivolts \\
\quad\code{samples} & The values, in those units; \code{null} where a sample is missing \\
\quad\code{fs} & The lead's sampling rate, when it differs \\
\code{meta} & Optional: the report \\
\quad\code{fields} & \code{[[label, text], \dots]}: the details; a field that names the patient is
  dropped \\
\quad\code{measurements} & \code{[\{name, value, unit, lead\}, \dots]}; \code{lead} is optional \\
\quad\code{statements} & \code{[\{source, program, by, time, lines\}, \dots]}: \code{source} is
  \code{computer}, \code{confirmed} or \code{edited}; \code{lines} is a list of texts; the others
  are optional \\
\quad\code{notes} & A list of texts \\
\end{xltabular}

A record must have at least one lead, a sampling rate for each, and at least one sample;
otherwise the page says which is missing (\code{this ECG has no leads}, \code{lead NAME: no
sampling rate}, \code{this ECG has no samples}).

\subsection{Pictures and PDF}
\label{sec:ecg-viewer-pictures}

A PNG or JPEG picture (a scanned or exported ECG) is shown as it is, and a PDF in the browser's
PDF viewer. The controls of the paper (layout, speed, gain, zoom and calipers) are grayed out,
because a picture's layout, speed and gain are fixed, and calipers need the signal.

% ==========================================================================================
\section{Opening ECGs}
\label{sec:ecg-viewer-open}

\begin{table}[htbp]
\centering
\caption{The ways to open ECGs in the ECG viewer.}
\label{tab:ecg-viewer-open}
\begin{tabularx}{\linewidth}{@{}P{4.2cm}L@{}}
\toprule
& \tabhead Opens \\
\midrule
\ui{Open record\dots} & One ECG: choose its file, or a WFDB record's \file{.hea} together with
  its signal files (and a multi-segment record's segments). \\
\ui{Open multiple records\dots} & A folder: every ECG in it and in its subfolders, one at a
  time. \\
Drop & The files dropped on the page, as with \ui{Open record\dots}. \\
\ui{Server folders\dots} & Shown only when the page is served by \cmd{muse-py serve}: the
  folders of the serving computer (\cref{sec:ecg-viewer-serve}). \\
\bottomrule
\end{tabularx}
\end{table}

\paragraph{Open record\dots} opens the browser's file chooser, which offers the files whose
extensions the viewer's readers claim: \file{.hea} and \file{.dat}, \file{.xml}, \file{.dcm} and
\file{.dicom}, \file{.edf} and \file{.bdf}, \file{.json}, the pictures (\file{.png},
\file{.jpg}, \file{.jpeg}, \file{.gif}, \file{.webp}) and \file{.pdf}. Several files may be
chosen at once. Of the files chosen, the header is the ECG shown and the others are its signal
files; without a header, the first file is shown.

\paragraph{Open multiple records\dots} opens the browser's folder chooser. The folder's ECGs are
its files, at every level, whose extension a reader claims (as above, except the signal files
\file{.dat}), in the order of their folder and then of their name. EDF files are listed with the
others. Hidden files and folders (names that start with \code{.}) are passed over, and the
segments and layout of a multi-segment record belong to it. A WFDB record's signal files are taken
from its own folder. The first ECG opens; then:
\begin{itemize}
\item \ui{\textless} and \ui{\textgreater} (or \key{$\leftarrow$} and \key{$\rightarrow$}) step to the
  previous and the next ECG, and the list between them picks one;
\item the name shows the ECG's path in the folder (without \file{.hea}) and its place:
  \code{ecg/2024-03/rest\_0003 (3 of 12)}.
\end{itemize}
A folder of more than 5000 ECGs is refused: \code{This folder holds more than 5000 ECGs: open a
subfolder.} A folder without one: \code{No ECG in this folder or its subfolders.} An ECG of the
folder that cannot be read says why when it is stepped to; the others still open.

\paragraph{Drop.} Files dropped anywhere on the page open as with \ui{Open record\dots}.

% ==========================================================================================
\section{The page}
\label{sec:ecg-viewer-page}

The header holds, from left to right: the name \ui{ECG}, the file's name and a line about the
record (its start, as UTC, the number of leads, the sampling rate and the length: \code{2024-02-20
09:20:00 $\cdot$ 12 leads $\cdot$ 500 Hz $\cdot$ 10.0 s}), the controls of the paper, and the
buttons to open files. The report and the paper fill the page below it.

\begin{xltabular}{\linewidth}{@{}P{4.4cm}L@{}}
\caption{The controls and keys of the ECG viewer.}
\label{tab:ecg-viewer-controls}\\
\toprule
{\tabhead Control} & {\tabhead Does} \\
\midrule
\endfirsthead
\toprule
{\tabhead Control} & {\tabhead Does} \\
\midrule
\endhead
\bottomrule
\endfoot
\ui{Layout} & \ui{3 $\times$ 4 + rhythm}, \ui{6 $\times$ 2 + rhythm}, \ui{12 $\times$ 1}
  (\cref{sec:ecg-viewer-paper}); only for a record with the twelve standard leads \\
\ui{Speed} & \ui{25 mm/s} or \ui{50 mm/s} \\
\ui{Gain} & \ui{5 mm/mV}, \ui{10 mm/mV} or \ui{20 mm/mV} \\
\ui{$-$} \ui{Fit} \ui{+} & Smaller; fit the width of the page; larger. Zoomed, the middle
  button shows the size, \ui{100 \%} being the paper's real size on a 96\,dpi screen; a click on
  it fits the width again \\
\ui{Calipers} & Measure between two points of the traces (\cref{sec:ecg-viewer-calipers}) \\
\key{Esc} & Clear the calipers' points \\
\ui{\textless}, the list, \ui{\textgreater}; \key{$\leftarrow$} \key{$\rightarrow$} & The previous
  ECG of the folder, any of them, the next (shown with \ui{Open multiple records\dots}) \\
\ui{Server folders\dots} & The folders of the computer serving the page (\cmd{muse-py serve}
  only) \\
\ui{Open record\dots} & One ECG (\cref{sec:ecg-viewer-open}) \\
\ui{Open multiple records\dots} & A folder of ECGs \\
The moon or sun button & Switches between the light and the dark theme \\
\end{xltabular}

The layout, the speed, the gain and the theme are remembered in the browser. For a picture or a
PDF, the controls of the paper are grayed out. When another program opens the page at a file
(\code{\#src=}), the buttons that open files, the folder's arrows and list, \ui{Server
folders\dots} and the theme button are hidden: the host decides what is shown
(\cref{sec:ecg-viewer-address}).

% ==========================================================================================
\section{The paper}
\label{sec:ecg-viewer-paper}

The ECG is drawn on standard paper: 25\,mm/s and 10\,mm/mV by default (or 50\,mm/s; 5 or
20\,mm/mV), a 1\,mm and 5\,mm grid, a 1\,mV calibration pulse (200\,ms wide) at the start of each
row, the lead's name at the start of each segment, and a mark where the lead changes. The foot of
the paper repeats the speed, the gain and the sampling rate (\code{25 mm/s 10 mm/mV 500 Hz}). The
strip is 10 seconds long, or the record's length when it is shorter.

A record with the twelve standard leads (I, II, III, aVR, aVL, aVF, V1 to V6, by these names or
close ones: \code{ECG II}, \code{Lead II}, \code{avr}, \code{MLII} for II) is laid out in one of
three layouts:

\begin{table}[htbp]
\centering
\caption{The layouts of a 12-lead ECG.}
\label{tab:ecg-viewer-layouts}
\begin{tabularx}{\linewidth}{@{}P{3.2cm}L@{}}
\toprule
\tabhead Layout & \tabhead Rows \\
\midrule
\ui{3 $\times$ 4 + rhythm} & Three rows of four columns of 2.5\,s: I, aVR, V1, V4; II, aVL, V2,
  V5; III, aVF, V3, V6; then lead II over the whole strip \\
\ui{6 $\times$ 2 + rhythm} & Six rows of two columns of 5\,s: I, V1; II, V2; III, V3; aVR, V4; aVL,
  V5; aVF, V6; then lead II over the whole strip \\
\ui{12 $\times$ 1} & Twelve rows, each lead over the whole strip, in the order I to V6 \\
\bottomrule
\end{tabularx}
\end{table}

Any other record (fewer leads, other names) shows each lead in full, one per row, and
\ui{Layout} is grayed out. The paper fits the width of the page; \ui{$-$} and \ui{+} step through
fixed sizes.

\subsection{Calipers}
\label{sec:ecg-viewer-calipers}

\ui{Calipers} turns the calipers on (the button is highlighted) or off. Click two points of the
paper: each snaps to the nearest sample of a trace, within 40\,ms of the click along the paper.
The viewer shows the interval between them, in milliseconds, and the voltage difference, in
microvolts, the second point's minus the first's: \code{412 ms $\cdot$ +150 $\mu$V}. A third click
starts again; \key{Esc} clears the points, and so does a change of layout, speed or gain.

\subsection{The report above the paper}
\label{sec:ecg-viewer-report}

Above the paper, the report gives what the file says about the ECG:
\begin{itemize}
\item its details: time, device, site, status, age, sex, filters, and a MUSE file's location and
  reason for the test (in a WFDB header, \code{Reason} is a statement);
\item its measurements: rates, PR, QRS, QT, QTc, axes, and, on a line per measurement, those
  given per lead (\code{ST I -0.07 $\cdot$ II -0.07 \dots{} mV}); millivolts with two decimals;
\item its statements, side by side: the computer's interpretation (\ui{Computer interpretation},
  with its program), and the physician's (\ui{Confirmed by}, with the name and the time, or
  \ui{Edited, not confirmed});
\item its notes.
\end{itemize}

\begin{phibox}
A field whose key names the patient anywhere in it (name, ID, MRN, accession or NHS number,
birth date, address, zip, phone, e-mail \dots) is never shown, even when the file holds it. Free
text (comment lines without a key, statements) is shown as written. The name of the physician
who confirmed the statements is shown as the file gives it.
\end{phibox}

\subsection{Light and dark}
\label{sec:ecg-viewer-themes}

The page has a light and a dark theme (\cref{fig:ecg-viewer-main,fig:ecg-viewer-dark}). Until a
theme is chosen, the page follows the system's setting. The moon button (in the light theme) or
the sun button (in the dark theme) switches the theme, and the choice is remembered in the
browser. \code{theme=light} or \code{theme=dark} in the address sets the theme for that page
without remembering it; Medorama passes its own page's theme this way.

% ==========================================================================================
\section{Running and deploying}
\label{sec:ecg-viewer-run}

\begin{xltabular}{\linewidth}{@{}P{5.0cm}L@{}}
\caption{The ways to run the ECG viewer.}
\label{tab:ecg-viewer-run}\\
\toprule
{\tabhead How} & {\tabhead What it opens} \\
\midrule
\endfirsthead
\toprule
{\tabhead How} & {\tabhead What it opens} \\
\midrule
\endhead
\bottomrule
\endfoot
The website, \url{https://medorama.org/viewers/ecg/} & The files you pick
  or drop, read in your browser: they never leave your computer \\
\file{viewers/ecg/index.html} from a clone, opened from disk (\code{file://}) & The same,
  offline \\
Any web server: copy the page's files (\file{index.html}, \file{style.css} and the \file{.js} files
  of \file{viewers/ecg}) & The same; and files on that server, by address (\code{\#src=URL}) \\
\cmd{muse-py serve} on the computer where the ECGs are & The same; and the folders of that
  computer, under a root you choose (\ui{Server folders\dots}) \\
Inside Medorama (\cmd{tools/install-medorama.sh}) & A patient's ECGs, as a tab of the patient's
  page, de-identified by Medorama's server (\cref{chap:viewers-in-medorama}) \\
\end{xltabular}

The page needs no build and no dependencies; it uses nothing but the browser. For a quick look
from a web server: \cmd{python3 -m http.server -{}-directory viewers/ecg}.

\subsection{Served from a computer: \texorpdfstring{\cmd{muse-py serve}}{muse-py serve}}
\label{sec:ecg-viewer-serve}

The Python package \file{viewers/ecg/python/muse} (\code{muse\_reader\_py}, the command
\cmd{muse-py}) serves the page together with the folders of the computer it runs on, so that the
ECGs of a server are read without copying them, through an SSH tunnel to its port (the commands
to install it on each operating system are at \url{https://medorama.org/docs/install.html}). \cmd{serve}
uses only the Python standard library; its module, \file{page\_server.py}, is the same file as
\cmd{dicom-py serve}'s.

\begin{lstlisting}[style=shell]
pip install ./viewers/ecg/python/muse         # copies the page into the package
muse-py serve                                 # opens http://127.0.0.1:8080/browse: your home folder
muse-py serve --root /data/ecg                # another folder
muse-py serve --set-password                  # once, on a shared machine (VitalStream's password)
\end{lstlisting}

\begin{lstlisting}[style=signature]
muse-py serve [--root DIR] [--host HOST] [--port N] [--auth-file FILE]
              [--set-password | --clear-password] [--no-auth] [--no-browser]
\end{lstlisting}

\begin{params}
\item[--root DIR] The folder served: only the files and folders under it are listed and served,
  every symbolic link resolved (a link that leads outside is refused, and not listed). Default:
  the home folder. It lies inside the root the package was installed with: a folder that
  contains that root is narrowed to it, one outside it is refused.
\item[--host HOST] Where it listens. Default: \code{127.0.0.1}. A request must name the server
  by a loopback name, an IP address or this host; any other name is refused (no DNS rebinding).
\item[--port N] Default: 8080.
\item[--auth-file FILE] The password file, shared with VitalStream and \cmd{dicom-py serve}.
  Default: \file{\textasciitilde/.config/vitalstream/auth.json}.
\item[--set-password] Set (or change) the password the page asks for, then exit
  (\code{New viewer password:}, at least 8 characters, then \code{Repeat it:}). For
  \cmd{muse-py serve} and \cmd{dicom-py serve}, a sign-in lasts 30 days, until the server stops,
  or until a browser signs out, which ends every sign-in; after a wrong password the next try,
  from any browser, waits 1\,s, then 2, 4 \dots{} up to 60\,s.
\item[--clear-password] Remove the password, then exit.
\item[--no-auth] Run without asking for the password.
\item[--no-browser] Do not open a browser window.
\end{params}

When it starts, \cmd{serve} prints its address (\code{ECG viewer: open http://127.0.0.1:8080/
(server folders: http://127.0.0.1:8080/browse)}), its root, and whether a password protects it;
then it opens the browser at \file{browse}, unless told not to.

\paragraph{The server folders.} \file{browse} (titled \ui{ECG viewer: server folders}) lists the
folders and the ECG files of a folder under the root: the files ending in \file{.hea},
\file{.xml}, \file{.json}, \file{.dcm}, \file{.pdf}, \file{.png}, \file{.jpg} and \file{.jpeg},
with their size; hidden entries are neither listed nor served, and a folder shows its first 5000
entries. A row \ui{..} leads to the folder above. A click on a file opens it in the viewer
(\code{\#src=files/PATH}); the browser's Back button returns to the list. The page opened at its
own address (\url{http://127.0.0.1:8080/}) shows \ui{Server folders\dots}, which leads to the list.

Files are sent as they are read, in part when a browser asks for a range (\code{Range}), under a
policy that runs nothing in them (\code{Content-Security-Policy: sandbox}); a kind of file the
viewers do not read is sent as a download. The sign-in cookie is signed with a key of the
running server too: browsers send a cookie to every port of a host, and another server on
127.0.0.1 cannot use it.

\paragraph{The requests.} \cmd{serve} answers \code{GET /} (the page), \code{GET browse?path=}
(a folder's records and subfolders, as HTML), \code{GET files/PATH} (a file), \code{GET
list/PATH/} (a folder's files as \code{\{"files": [\dots]\}}, the PACS viewer's \code{\#list=}; and
\code{list/PATH} without a final \code{/}, the file itself) and \code{GET served}, which answers
200 \code{\{"viewer", "root"\}}: when it does, the page shows its \ui{Server folders\dots} button.

\paragraph{Fixed at install time.} Two options are fixed when the package is installed, from the
environment:

\begin{table}[htbp]
\centering
\caption{The install-time variables of \cmd{muse-py} (and \cmd{dicom-py}).}
\label{tab:ecg-viewer-serve-build}
\begin{tabularx}{\linewidth}{@{}P{3.0cm}P{2.8cm}L@{}}
\toprule
\tabhead Variable & \tabhead Default & \tabhead Effect \\
\midrule
\code{VIEWER\_ROOT} & The whole file system & The most \cmd{serve} ever serves;
  \code{\textasciitilde} is the home folder of whoever runs it. A \opt{-{}-root} that contains it
  is narrowed to it, and one outside it is refused (\code{\dots{} is not a folder inside this
  viewer's root}): then give a \opt{-{}-root} inside it \\
\code{VIEWER\_PASSWORD} & \code{1} & \code{0}: without password support (\opt{-{}-set-password}
  fails, the page never asks for one) \\
\bottomrule
\end{tabularx}
\end{table}

\begin{lstlisting}[style=shell]
VIEWER_ROOT=/data/ecg VIEWER_PASSWORD=1 pip install ./viewers/ecg/python/muse
\end{lstlisting}

An editable install (\cmd{pip install -e}, as \cmd{tools/install-medorama.sh} makes) serves the
page from the clone instead of a copy. The two variables still apply when they are given at
install time (they are written into the clone, in a file git ignores); without them, its root is
the whole file system, with password support.

\paragraph{On a server where the ECGs are:} install with a \code{VIEWER\_ROOT} that holds them,
set a password once (\cmd{muse-py serve -{}-set-password}), and start it without a browser, in
\cmd{tmux} or as a service: \cmd{muse-py serve -{}-root /data/ecg -{}-no-browser}. From your
computer: \cmd{ssh -N -L 8080:127.0.0.1:8080 user@server}, then open
\url{http://127.0.0.1:8080/}. Without a password, anyone who can connect to the port can read the
files under the root. The security model shared by the viewers is described in
\cref{chap:viewers-alone}.

% ==========================================================================================
\section{muse-py: MUSE files outside the page}
\label{sec:ecg-viewer-muse-py}

\cmd{muse-py} reads GE MUSE resting ECG XML files with no dependencies, as the page does: to the
JSON form, to a WFDB record whose header comments carry the report, or to an ECG database to
import into Medorama.

\begin{lstlisting}[style=shell]
pip install -e viewers/ecg/python/muse
muse-py json ECG.xml > ecg.json          # the JSON form, with its report (meta)
muse-py wfdb ECG.xml out/                # a WFDB record (format 16, 1 uV per unit)
muse-py index /path/to/muse -o ecg.csv   # an ECG database: one row per file
muse-py --version
\end{lstlisting}

\apientry{muse-py json}{The ECG in the JSON form}
\begin{lstlisting}[style=signature]
muse-py json FILE
\end{lstlisting}
Writes the ECG of the MUSE file \file{FILE} in the viewer's JSON form, with its report in
\code{meta}, on the standard output (one line).

\apientry{muse-py wfdb}{The ECG as a WFDB record}
\begin{lstlisting}[style=signature]
muse-py wfdb FILE OUT_DIR [--name NAME]
\end{lstlisting}
Writes \file{NAME.hea} and \file{NAME.dat} into \file{OUT\_DIR} (created if needed), and prints
the header's path: the leads in format 16, at 1000 units per millivolt (1\,$\mu$V per unit), with
the details, the measurements and the statements in the header's comments (\code{\# Dx: \dots},
\code{\# Dx confirmed: \dots}, \code{\# Confirmed by: \dots}), each value on one line (line breaks in
it become spaces), which the ECG viewer reads back as the report.
\begin{params}
\item[--name NAME] The record's name. Default: the file's name without its extension.
\end{params}

\apientry{muse-py index}{An ECG database}
\begin{lstlisting}[style=signature]
muse-py index FOLDER -o CSV
\end{lstlisting}
Reads every MUSE file (\file{.xml}) under \file{FOLDER}, passing over hidden files and folders,
and writes a CSV, one row per file, to
import into Medorama (\cref{chap:ecg-imaging}). It prints \code{CSV: N ECG(s)}, and
\code{; M file(s) not read} with each of them on the standard error.
\begin{params}
\item[-o, --output CSV] The CSV to write (required); the file paths in it are relative to its
  folder.
\end{params}

The columns of the database:
\begin{itemize}
\item \code{mrn} (the MUSE \code{PatientDemographics/PatientID}, in a column named \code{mrn}:
  Medorama links each row to its patient by it, as an MRN), \code{time}, \code{category}
  (\code{ecg}), \code{name} (\code{12-lead ECG}), and \code{value} and \code{unit} (the
  ventricular rate);
\item one column per measurement: \code{PR (ms)}, \code{QT (ms)}, \code{QTc (ms)} and the others;
\item \code{interpretation} (12SL) and \code{confirmed};
\item \code{file}: its path, relative to the CSV.
\end{itemize}

\begin{phibox}
No patient identity is written into the JSON or WFDB outputs. The database holds the patient ID
(\code{mrn}\phimark), the key that links each ECG to its patient: keep it as PHI, with Medorama's
databases, never in the repository.
\end{phibox}

% ==========================================================================================
\section{Addresses}
\label{sec:ecg-viewer-address}

\code{index.html\#src=URL} shows the file at \code{URL}; a WFDB header's signal files, and a
multi-segment record's segments, are read beside it (\code{URL}'s folder followed by their
names). The file is fetched with the page's own credentials. Changing the fragment shows another
file without reloading the page.

\begin{table}[htbp]
\centering
\caption{The fragments of the ECG viewer's address.}
\label{tab:ecg-viewer-address}
\begin{tabularx}{\linewidth}{@{}P{4.6cm}L@{}}
\toprule
\tabhead Fragment & \tabhead Means \\
\midrule
\code{src=URL} & The file to show \\
\code{name=\dots} & The name to show \\
\code{theme=light} or \code{dark} & The theme \\
\code{compact=1} & The paper alone, fitted to the frame (shown inside another page): no header,
  a shorter report \\
\code{shift=SECONDS} & For a host that de-identifies in the page: the report's dates, and the
  record's start in the header line, move back by it \\
\code{cap90=1} & The same: ages above 89 years shown as 90+ (an \code{Age} or \code{<age>}
  field; in days, weeks or months when its unit says so) \\
\code{who=\dots}, \code{what=\dots}, \code{pos=\dots}, \code{nav=\dots} & A host's context, shown in
  place of the file's name (\cref{chap:viewers-alone}) \\
\bottomrule
\end{tabularx}
\end{table}

Medorama does not need \code{shift} and \code{cap90}: its server de-identifies the files
themselves. Example, from a web server that serves the page and a folder \file{ecgs/}:
\code{index.html\#src=ecgs/sinus\_st\_depression.hea\&theme=light}.

% ==========================================================================================
\section{The examples}
\label{sec:ecg-viewer-examples}

\file{viewers/ecg/examples/} holds synthetic ECGs to try the viewer with: every value is
invented, and no file describes a real person. They are written by
\file{tools/make-viewer-examples.py} with the demo's generators. On the website, the examples page
(\url{https://medorama.org/viewers/ecg/examples/}) has a link that opens
each file in the live viewer; from a clone, open \file{viewers/ecg/index.html} and choose a file
with \ui{Open record\dots} (a WFDB record: its \file{.hea} and \file{.dat} together), or the whole
folder with \ui{Open multiple records\dots}.

\begin{table}[htbp]
\centering
\caption{The ECG viewer's examples.}
\label{tab:ecg-viewer-examples}
\begin{tabularx}{\linewidth}{@{}P{5.0cm}L@{}}
\toprule
\tabhead File & \tabhead What \\
\midrule
\file{sinus\_st\_depression.hea}, \file{.dat} & A WFDB record of twelve leads (500\,Hz, 10\,s), sinus
  rhythm with inferolateral ST depression; the cart's report (measurements, the computer's
  statements) in the header's comments \\
\file{atrial\_fibrillation\_muse.xml} & A GE MUSE resting ECG XML file: atrial fibrillation with a
  rapid ventricular response, the measurements and the 12SL statements \\
\file{normal\_dicom.dcm} & A DICOM 12-lead ECG waveform object: normal sinus rhythm, with its
  annotations as the report \\
\file{normal\_printout.pdf} & A PDF printout of a 12-lead ECG, shown as it is \\
\file{normal\_printout.png} & The same as a picture \\
\bottomrule
\end{tabularx}
\end{table}

% ==========================================================================================
\section{Files and tests}
\label{sec:ecg-viewer-files}

\begin{table}[htbp]
\centering
\caption{The files of the ECG viewer (\file{viewers/ecg}).}
\label{tab:ecg-viewer-files}
\begin{tabularx}{\linewidth}{@{}P{3.6cm}L@{}}
\toprule
\tabhead File & \tabhead Contents \\
\midrule
\file{readers.js} & The registry of readers: one reader per format, each giving the same record.
  Its header comment documents the record, the JSON form, the reader interface and how to add a
  format (\cref{sec:readers-ecg}). \\
\file{ecg.js} & The record model and its helpers: the JSON form (\code{ECG.readJson}), units to
  millivolts, metadata from ``key: value'' lines, the report, the standard lead names, and
  \code{ECG.checked}, which refuses a record that the paper cannot draw. \\
\file{wfdb.js} & The WFDB reader: headers and signal files, single- and multi-segment. \\
\file{muse.js} & The GE MUSE XML reader, with its own small XML parser. \\
\file{dicom-ecg.js} & The DICOM ECG reader, over \file{dicom.js}. \\
\file{dicom.js} & The DICOM parser, a copy of the PACS viewer's. \\
\file{edf.js} & The EDF, EDF+, BDF and BDF+ reader. \\
\file{paper.js} & The paper drawing (SVG in millimeters) and the layouts. \\
\file{host.js} & A host's context (\code{who}, \code{what}, \code{pos}, \code{nav}), shown in
  place of the file's name; the same file in each viewer. \\
\file{app.js}, \file{index.html}, \file{style.css} & The page. \file{index.html} loads the
  readers after \file{readers.js}, in the order in which the page lists the formats. \\
\file{python/muse/} & \cmd{muse-py}: the MUSE reader in Python (\cmd{json}, \cmd{wfdb},
  \cmd{index}) and \cmd{serve}. \\
\file{examples/} & Synthetic ECGs to try the viewer with. \\
\file{tests/} & The tests of the page (\file{readers.test.js}, \file{ecg.test.js},
  \file{muse.test.js}, \file{dicom.test.js}, \file{edf.test.js}, \file{host.test.js}) and of
  \cmd{muse-py}. \\
\bottomrule
\end{tabularx}
\end{table}

\begin{lstlisting}[style=shell]
node --test viewers/ecg/tests/*.test.js        # the page's readers and drawing
pytest viewers/ecg/python/muse/tests           # muse-py, serve (its root and its password)
\end{lstlisting}

The MUSE readers are tested on \file{viewers/ecg/tests/fixtures/synthetic\_muse.xml}, made by
\file{make\_muse.py} with an invented patient and signals; \file{synthetic\_muse.json} is the
Python reader's output, which the JavaScript reader's tests compare against. The reader was
written from the format, independently of GE's GPL \file{musexmlex.py}.
